
Can fine-tuning teach a model to perform tasks its architecture fundamentally cannot support? This work demonstrates that parameter-efficient fine-tuning (PEFT) methods adapt existing architectural capabilities but cannot add missing ones. When pretrained models achieve below-random zero-shot accuracy on a task, architectural incompatibility prevents successful adaptation regardless of fine-tuning strategy.

We evaluate Mamba-130M, a causal state-space model, on GLUE classification tasks requiring different reasoning capabilities. On sentiment classification (SST-2), the model achieves 81\% zero-shot accuracy (31 percentage points above random, p < 0.001), demonstrating language understanding compatible with unidirectional processing. However, on paraphrase detection (QQP), the same checkpoint achieves 38\% accuracy---12 percentage points below the 50\% random baseline (p = 0.003). This below-random performance reveals architectural incompatibility: symmetric question comparison requires bidirectional reasoning, which causal state-space architectures cannot provide.

These results establish zero-shot evaluation as a mandatory gate before PEFT investment. Below-random accuracy signals architecturally impossible tasks where fine-tuning fails. This principle challenges assumptions underlying PEFT transfer across architectures: methods developed for transformers do not automatically transfer to sub-quadratic models. As efficient architectures proliferate, validating architectural alignment becomes essential to prevent wasted compute on structurally impossible experiments.
